% -----------------------------------------------------------------------------
% Drop-in subsection: Unified defect action
% Suggested insertion: Section ``Theory: generative bandwidth relativity'',
% immediately after Eq. \eqref{eq:rest-lapse} and before the paragraph beginning
% ``The spatial source law is distinct.''
% -----------------------------------------------------------------------------
% Optional preamble helpers, if you want named commands:
% \newcommand{\Xiang}{\Xi}
% \newcommand{\Urad}{\Upsilon}

\subsection{Unified defect action: lapse and holonomy as projections}\label{sec:unified-defect-action}

Equations~\eqref{eq:unified}--\eqref{eq:rest-lapse} describe the
\emph{temporal} projection of a defect: a persistent unresolved region removes
ordinary update capacity from the active frontier, producing a lapse
$\alpha=1-\Omega$.  The spatial experiments show that this lapse projection is
not, by itself, identical to spatial holonomy.  A pure $1\to3$ refinement sink can
produce a useful-update lapse while failing to reduce the refinement-weighted
circumference of exterior annuli.  Conversely, an explicit sector or linkage
change is visible as holonomy in the $q$-sector positive control.  The more
natural unified object is therefore not a single scalar curvature variable, but
a local \emph{generative defect action} whose distinct projections are observed
as lapse, radial throat response, and angular holonomy.

Let $\mathcal{D}_m(R)$ denote the defect action accumulated inside a weighted
radius $R$.  Operationally, $\mathcal{D}_m$ is the local count of distinction
operations and linkage opportunities made unavailable to ordinary vacuum
propagation by a persistent defect.  We write three projections:
\begin{align}
  \Omega(R) &= P_\tau[\mathcal{D}_m(R)] && \text{temporal lapse / update starvation},\\
  \Upsilon(R) &= P_R[\mathcal{D}_m(R)] && \text{radial throat or interior stretch},\\
  \Xi(R) &= P_\theta[\mathcal{D}_m(R)]
     =1-\frac{C_m(R)}{C_0(R)} && \text{angular holonomy / circumference pinch}.
\end{align}
Here $C_0(R)$ is the matched vacuum circumference and $C_m(R)$ is the measured
circumference in the defect run, using the same refinement-weighted metric.  In
a perfect $q$-sector disclination,
\begin{equation}
  \Xi = \frac{6-q}{6},
  \qquad
  \frac{C_q}{C_6}=\frac{q}{6},
\end{equation}
which is the discrete positive-control version of the conical holonomy of
$2+1$-dimensional gravity~\cite{DeserJackiwTHooft1984}.

The corresponding effective axisymmetric $2+1$D line element can be written as
\begin{equation}\label{eq:deu-projected-metric}
  ds_{\rm DEU}^2
  =-[1-\Omega(R)]^2dt^2
  +[1+\Upsilon(R)]^2dR^2
  +[1-\Xi(R)]^2R^2d\theta^2 .
\end{equation}
For a moving probe this gives
\begin{equation}\label{eq:deu-probe-propertime}
  \frac{d\tau}{dt}
  =\sqrt{
     [1-\Omega(R)]^2
     -[1+\Upsilon(R)]^2\dot R^{\,2}
     -[1-\Xi(R)]^2 R^2 \dot\theta^{\,2}}
  .
\end{equation}
In the no-holonomy, locally flat limit $\Upsilon=\Xi=0$ and
$v^2=\dot R^{\,2}+R^2\dot\theta^{\,2}$, Eq.~\eqref{eq:deu-probe-propertime}
reduces to the bandwidth clock law
\begin{equation}
  \frac{d\tau}{dt}=\sqrt{(1-\Omega)^2-v^2}.
\end{equation}
Thus the lapse equation is the temporal projection of the same defect action
whose angular projection is measured by $\Xi$.

The simulations suggest that $\Omega$ and $\Xi$ should not be equated as a
pointwise identity.  Instead, $\Omega$ is a comparatively continuous bandwidth
projection, while $\Xi$ is a quantized or thresholded holonomy projection.  A
useful phenomenological form is
\begin{equation}\label{eq:stitch-threshold}
  \Xi(\Omega)
  =\frac{1}{6}
   \left\lfloor
     \frac{\Omega-\Omega_c}{\Delta\Omega_{\rm stitch}}
   \right\rfloor_+,
\end{equation}
where $\Omega_c$ is the topological slack that the local foam can absorb before
an exterior sector is lost, and $\Delta\Omega_{\rm stitch}$ is the additional
defect action required to drop another effective sector.  In a smooth weak-field
regime one may observe an approximate proportionality $\Xi\simeq\kappa\Omega$,
but in the native discrete regime the holonomy channel is expected to exhibit
stair-step behavior.

This formulation reconciles the two main Layer~3 findings.  A refinement sink is
a Zeno sponge: it has a lapse projection $\Omega>0$ but can leave $\Xi\simeq0$ in
the exterior refinement-weighted metric.  A topology stitch activates the angular
projection $\Xi>0$.  When $\Xi>0$ occurs together with radial/interior response
$\Upsilon>0$, the resulting geometry is not a pure cone but a throat: the
exterior circumference is pinched while the interior weighted area can remain
large or even inflate.
